\section{Zassenhaus factorization: every exponent, three recursions, a tree formula, and convergence}\label{sec:zassenhaus}
BCH combines a product of exponentials into one exponential. Zassenhaus
factorization performs the ordered operation in the opposite direction:
it writes the exponential of a sum as an ordered product of exponentials of
homogeneous Lie corrections. The order of the factors is part of the
statement. Recursive constructions are central in the computational
literature \cite{casasmuruanadinic}; explicit reorganizations are
considered in \cite{kimura}. We prove existence and uniqueness, three
recursive constructions, a finite nonrecursive formula for every exponent,
a globally convergent ordered-integral expansion, and a quantitative
convergence theorem for the product.

\subsection{Existence, uniqueness, and the residual algorithm}
\begin{theorem}[Formal Zassenhaus factorization]\label{thm:zassenhaus}
There is a unique sequence of homogeneous Lie polynomials $C_n(X,Y)$ of
degree $n$, $n\geq2$, such that
\begin{equation}\label{eq:zassenhaus}
 \boxed{\displaystyle
 e^{t(X+Y)}=e^{tX}e^{tY}\prod_{n=2}^{\infty}e^{t^nC_n(X,Y)}
 =e^{tX}e^{tY}e^{t^2C_2}e^{t^3C_3}e^{t^4C_4}\cdots.}
\end{equation}
The infinite product is ordered by increasing $n$ from left to right and is
well defined coefficientwise in $t$: modulo $t^{N+1}$ only the factors
through $C_N$ matter. Define
\begin{equation}\label{eq:R1}
 R_1(t)=e^{-tY}e^{-tX}e^{t(X+Y)},
\end{equation}
and recursively, for $n\geq2$,
\begin{equation}\label{eq:residual}
 C_n=\coeff{t^n}R_{n-1}(t)=\coeff{t^n}\log R_{n-1}(t),\qquad
 R_n(t)=e^{-t^nC_n}R_{n-1}(t).
\end{equation}
Then $R_n(t)=I+O(t^{n+1})$ for every $n\geq1$, and \eqref{eq:residual}
computes all the exponents.
\end{theorem}
\begin{proof}
Work in $\widehat T[[t]]$, or equivalently in $\widehat T$ with $t$ tracking
total degree. The coefficient of $t^1$ in $R_1$ is $-Y-X+(X+Y)=0$, so
$R_1=I+O(t^2)$. Suppose inductively that $R_{n-1}=I+t^nC_n+O(t^{n+1})$.
Then $\log R_{n-1}=t^nC_n+O(t^{n+1})$, which gives the second expression
for $C_n$ in \eqref{eq:residual}, and
$R_n=e^{-t^nC_n}R_{n-1}=(I-t^nC_n+O(t^{2n}))(I+t^nC_n+O(t^{n+1}))
=I+O(t^{n+1})$. Thus the recursion removes one degree at a time.

Each residual is group-like in the completed Hopf algebra of
\cref{sec:hopf} (extended by the central scalar $t$): $R_1$ is a product of
exponentials of the primitive elements $-tY$, $-tX$, $t(X+Y)$, hence
group-like by \cref{lem:grouplike}; and if $R_{n-1}$ is group-like and
$C_n$ is a Lie polynomial, then $e^{-t^nC_n}$ is group-like and so is the
product $R_n$. To see that $C_n$ is a Lie polynomial, compare the
coefficients of $t^n$ in $\Delta R_{n-1}=R_{n-1}\otimes R_{n-1}$: on the
left it is $\Delta C_n$; on the right, since
$R_{n-1}=I+t^nC_n+O(t^{n+1})$ has no terms of $t$-degree between $1$ and
$n-1$, it is $C_n\otimes1+1\otimes C_n$. Hence $C_n$ is primitive, and by
\cref{thm:friedrichs} it is a Lie polynomial; it is homogeneous of degree
$n$ in the letters because every input has matching degree in $t$ and in
$X,Y$. This completes the induction.

Undoing the residuals, $R_1=e^{t^2C_2}e^{t^3C_3}\cdots e^{t^NC_N}R_N$ with
$R_N=I+O(t^{N+1})$, so modulo $t^{N+1}$ the product of the first $N-1$
factors equals $R_1$; multiplying by $e^{tX}e^{tY}$ on the left proves
\eqref{eq:zassenhaus} degree by degree. For uniqueness, suppose
$e^{t(X+Y)}=e^{tX}e^{tY}\prod_ne^{t^nC_n'}$ with homogeneous $C_n'$ of
degree $n$. Then $\prod_{n\geq2}e^{t^nC_n'}=R_1$. If $C_j'=C_j$ for
$j<n$, then multiplying on the left by the inverses of the first factors
gives $\prod_{j\geq n}e^{t^jC_j'}=R_{n-1}$, and comparing coefficients of
$t^n$ gives $C_n'=C_n$. Induction proves uniqueness.
\end{proof}
When the identity is written as $e^{-tX}e^{t(X+Y)}=\prod_{n\geq1}e^{t^nC_n}$,
as on the source page, the convention \emph{must} be $C_1=Y$; the notation
would otherwise omit a necessary first-degree factor. Extracting $t^n$
from $R_{n-1}=e^{-t^{n-1}C_{n-1}}\cdots e^{-t^2C_2}R_1$ gives the finite
coefficient recurrence
\begin{equation}\label{eq:zassfinitecoeff}
 C_n=\sum_{\substack{k_2,\ldots,k_{n-1},a,b,c\geq0\\
                 \sum_{j=2}^{n-1}jk_j+a+b+c=n}}
 \Bigl(\prod_{j=n-1}^{2}\frac{(-C_j)^{k_j}}{k_j!}\Bigr)
 \frac{(-Y)^a}{a!}\frac{(-X)^b}{b!}\frac{(X+Y)^c}{c!},
\end{equation}
in which the product over $j$ is taken in \emph{descending} order of $j$
from left to right. It involves associative products, but the result is a
Lie polynomial by the theorem, and \eqref{eq:dsw} converts it into nested
commutators.

\subsection{A recursion consisting entirely of commutators}
Define the right logarithmic derivatives $F_n(t)=R_n'(t)R_n(t)^{-1}$.
\begin{proposition}[Differential Zassenhaus algorithm]\label{prop:zassdiff}
\begin{align}
 F_1(t)&=e^{-t\ad_Y}\bigl(e^{-t\ad_X}-I\bigr)Y
 =\sum_{k\geq1}t^kf_{1,k},\qquad
 f_{1,k}=(-1)^k\sum_{j=1}^{k}\frac{\ad_Y^{k-j}\ad_X^{j}Y}{(k-j)!\,j!},
 \label{eq:F1}\\
 C_n&=\frac1n\coeff{t^{n-1}}F_{n-1}(t),\qquad
 F_n(t)=e^{-t^n\ad_{C_n}}\bigl(F_{n-1}(t)-nt^{n-1}C_n\bigr)\quad(n\geq2).
 \label{eq:Frec}
\end{align}
Writing $F_n=\sum_{k\geq0}t^kf_{n,k}$, with $f_{n,k}=0$ for $k<0$, the
coefficients satisfy the fully finite update
\begin{equation}\label{eq:frec}
 \boxed{\displaystyle
 f_{n,k}=\sum_{j=0}^{\lfloor k/n\rfloor}\frac{(-1)^j}{j!}
 \ad_{C_n}^{\,j}f_{n-1,k-nj}-nC_n\,\delta_{k,n-1},\qquad
 C_n=\frac{f_{n-1,n-1}}{n}.}
\end{equation}
Only previously computed homogeneous Lie polynomials occur; this is a
Lie-only algorithm at every degree.
\end{proposition}
\begin{proof}
For a product $ABC$ of differentiable invertible curves,
$(ABC)'(ABC)^{-1}=A'A^{-1}+A(B'B^{-1})A^{-1}+AB(C'C^{-1})B^{-1}A^{-1}$.
With $A=e^{-tY}$, $B=e^{-tX}$, $C=e^{t(X+Y)}$ we have $A'A^{-1}=-Y$,
$B'B^{-1}=-X$, $C'C^{-1}=X+Y$, and conjugations are adjoint exponentials
by \cref{thm:campbell}:
\begin{align*}
 F_1&=-Y-e^{-t\ad_Y}X+e^{-t\ad_Y}e^{-t\ad_X}(X+Y)\\
 &=-Y-e^{-t\ad_Y}X+e^{-t\ad_Y}X+e^{-t\ad_Y}e^{-t\ad_X}Y
 =e^{-t\ad_Y}\bigl(e^{-t\ad_X}-I\bigr)Y,
\end{align*}
using $e^{-t\ad_X}X=X$ and $e^{-t\ad_Y}Y=Y$. Expanding both adjoint
exponentials gives the double series
$\sum_{i\geq0,j\geq1}(-1)^{i+j}t^{i+j}\ad_Y^i\ad_X^jY/(i!j!)$, and
collecting $t^k$ gives $f_{1,k}$.

Since $R_{n-1}=I+t^nC_n+O(t^{n+1})$, we have
$R_{n-1}'=nt^{n-1}C_n+O(t^n)$ and $R_{n-1}^{-1}=I+O(t^n)$, so
$F_{n-1}=nt^{n-1}C_n+O(t^n)$; in particular $f_{n-1,k}=0$ for $k<n-1$
and $f_{n-1,n-1}=nC_n$, the first part of \eqref{eq:Frec}. Differentiating
$R_n=e^{-t^nC_n}R_{n-1}$,
\[
 F_n=\Bigl(\frac{\dd}{\dd t}e^{-t^nC_n}\Bigr)e^{t^nC_n}
 +e^{-t^nC_n}F_{n-1}e^{t^nC_n}
 =-nt^{n-1}C_n+e^{-t^n\ad_{C_n}}F_{n-1},
\]
because $t^nC_n$ commutes with its derivative and conjugation by
$e^{-t^nC_n}$ is $e^{-t^n\ad_{C_n}}$. As $e^{-t^n\ad_{C_n}}C_n=C_n$, the
correction may be moved inside the operator exponential, giving the second
part of \eqref{eq:Frec}. Expanding
$e^{-t^n\ad_{C_n}}=\sum_j(-1)^jt^{nj}\ad_{C_n}^j/j!$ and extracting the
coefficient of $t^k$ gives \eqref{eq:frec}; the term $-nt^{n-1}C_n$ is
annihilated by $\ad_{C_n}^j$ for $j\geq1$ and contributes only through
$j=0$.
\end{proof}

\subsection{The exponents displayed on the source page}
\begin{proposition}\label{prop:zassdisplayed}
\begin{align}
 C_2&=-\tfrac12[X,Y],\label{eq:c2}\\
 C_3&=\tfrac16[X,[X,Y]]+\tfrac13[Y,[X,Y]]
     =\tfrac16[X,[X,Y]]-\tfrac13[[X,Y],Y],\label{eq:c3}\\
 C_4&=-\tfrac1{24}[X,[X,[X,Y]]]-\tfrac18[Y,[X,[X,Y]]]-\tfrac18[Y,[Y,[X,Y]]]
 \label{eq:c4}\\
    &=-\tfrac1{24}\bigl([[[X,Y],X],X]+3[[[X,Y],X],Y]+3[[[X,Y],Y],Y]\bigr).
 \label{eq:c4left}
\end{align}
\end{proposition}
\begin{proof}
From \eqref{eq:F1},
\begin{align*}
 f_{1,1}&=-[X,Y],\\
 f_{1,2}&=\ad_Y\ad_XY+\tfrac12\ad_X^2Y=[Y,[X,Y]]+\tfrac12[X,[X,Y]],\\
 f_{1,3}&=-\tfrac12\ad_Y^2\ad_XY-\tfrac12\ad_Y\ad_X^2Y-\tfrac16\ad_X^3Y.
\end{align*}
By \eqref{eq:frec}, $C_2=f_{1,1}/2=-\tfrac12[X,Y]$. Next,
$f_{2,2}=f_{1,2}-\ad_{C_2}f_{1,0}=f_{1,2}$ since $f_{1,0}=0$, so
$C_3=f_{2,2}/3$, which is \eqref{eq:c3}; the second form uses
$[Y,[X,Y]]=-[[X,Y],Y]$. Then
$f_{2,3}=f_{1,3}-\ad_{C_2}f_{1,1}=f_{1,3}+\tfrac12[[X,Y],[X,Y]]=f_{1,3}$,
and $f_{3,3}=f_{2,3}-\ad_{C_3}f_{2,0}=f_{2,3}$, so $C_4=f_{1,3}/4$, which is
\eqref{eq:c4}. For \eqref{eq:c4left}, antisymmetry gives
$[[X,Y],X]=-[X,[X,Y]]$, hence $[[[X,Y],X],X]=[X,[X,[X,Y]]]$,
$[[[X,Y],X],Y]=[Y,[X,[X,Y]]]$, and $[[[X,Y],Y],Y]=[Y,[Y,[X,Y]]]$.
Substituting these into \eqref{eq:c4} gives \eqref{eq:c4left}, exactly the
left-nested form on the source page. Since the corrections at each step
were computed exactly, no higher-order approximation has been left implicit.
\end{proof}
The complete fifth and sixth exponents are printed in \cref{app:tables}
(\cref{tab:zass56}), and all exponents through degree twelve accompany the
article as data. There is also a left-oriented companion factorization:
replacing $t$ by $-t$ in \eqref{eq:zassenhaus} and inverting both sides
reverses the factor order and gives
\begin{equation}\label{eq:zassleft}
 e^{t(X+Y)}=\cdots e^{t^4\widehat C_4}e^{t^3\widehat C_3}e^{t^2\widehat C_2}e^{tY}e^{tX},
 \qquad \widehat C_n=(-1)^{n+1}C_n,
\end{equation}
since $(e^{(-t)^nC_n})^{-1}=e^{(-1)^{n+1}t^nC_n}$.

\subsection{A finite nonrecursive formula using weighted trees}
The previous algorithms determine the exponents recursively. For a finite
nonrecursive expression at every fixed order, first define the explicit
associative polynomials
\begin{equation}\label{eq:An}
 A_n=\coeff{t^n}R_1(t)=\sum_{a+b+c=n}\frac{(-1)^{a+b}}{a!\,b!\,c!}\,Y^aX^b(X+Y)^c
 \qquad(n\geq0),
\end{equation}
where $(X+Y)^c$ is the sum of all $2^c$ words of length $c$; $A_0=I$ and
$A_1=0$. Coefficient extraction from the ordered product
$R_1=\prod_{j\geq2}e^{t^jC_j}$ gives the forward partition recurrence
\begin{equation}\label{eq:Cpartitions}
 C_n=A_n-\sum_{\substack{k_2,\ldots,k_{n-1}\geq0\\\sum_{j=2}^{n-1}jk_j=n}}
 \frac{C_2^{k_2}C_3^{k_3}\cdots C_{n-1}^{k_{n-1}}}{k_2!\,k_3!\cdots k_{n-1}!},
\end{equation}
with the factors in \emph{increasing} order of degree. Indeed, choosing
from each factor $e^{t^jC_j}$ the term $t^{jk_j}C_j^{k_j}/k_j!$ with
$\sum jk_j=n$ gives all contributions to $t^n$; the unique choice with
$k_n=1$ contributes $C_n$, and every other choice has $k_n=0$ and at
least two factors. The sum is empty for $n=2,3$.

\begin{definition}\label{def:trees}
Let $\cT_n$ be the set of finite plane rooted trees with a positive integer
weight at each vertex, subject to: the root has weight $n$; every weight is
at least $2$; and each internal vertex $v$ of weight $n_v$ has $r\geq2$
ordered children whose weights $2\leq j_1\leq j_2\leq\cdots\leq j_r$ are
weakly increasing from left to right and sum to $n_v$. Leaves carry no
further restriction. Let $I(T)$ be the number of internal vertices, let
$m_{v,j}$ be the number of children of weight $j$ at the internal vertex
$v$, and let $\ell_1,\ldots,\ell_q$ be the weights of the leaves read from
left to right in planar order. Define
\begin{equation}\label{eq:treeweight}
 \omega(T)=\frac{(-1)^{I(T)}}{\prod_{v\ \mathrm{internal}}\prod_{j\geq2}m_{v,j}!},
 \qquad
 A(T)=A_{\ell_1}A_{\ell_2}\cdots A_{\ell_q}.
\end{equation}
Children of equal weight are distinct ordered slots whose subtrees may
differ; interchanging two unequal subtrees changes $A(T)$ in general,
because the $A_j$ do not commute.
\end{definition}

\begin{theorem}[All-terms Zassenhaus tree formula]\label{thm:tree}
For every $n\geq2$, $\cT_n$ is finite and
\begin{equation}\label{eq:treeformula}
 \boxed{\displaystyle C_n=\sum_{T\in\cT_n}\omega(T)\,A(T).}
\end{equation}
This is a finite formula involving only factorials, sums, and products of
$X$ and $Y$. An explicitly nested-commutator version is
\begin{equation}\label{eq:treelie}
 C_n=\frac1n\sum_{T\in\cT_n}\omega(T)\sum_{w\in\{X,Y\}^n}\coeff{w}A(T)\,R(w).
\end{equation}
\end{theorem}
\begin{proof}
Finiteness: the leaf weights are at least $2$ and sum to $n$ (every
internal vertex's weight is the sum of its children's weights, so by
induction the root weight is the sum of the leaf weights), so there are at
most $\lfloor n/2\rfloor$ leaves; since every internal vertex has at least
two children, a plane tree with $q$ leaves has fewer than $q$ internal
vertices; and each vertex weight is an integer in $[2,n]$. Hence only
finitely many shapes and weight assignments occur.

We prove \eqref{eq:treeformula} by strong induction on $n$. For $n=2,3$
the sum in \eqref{eq:Cpartitions} is empty, $C_n=A_n$, and $\cT_n$ consists
of the single-vertex tree (a root of weight $n$ with no children, since
no weakly increasing list of at least two integers $\geq2$ sums to $2$ or
$3$), for which $\omega=1$ and $A(T)=A_n$. For $n\geq4$, assume the formula
for all degrees below $n$. In \eqref{eq:Cpartitions}, a choice of exponents
$(k_2,\ldots,k_{n-1})$ with $\sum jk_j=n$ corresponds bijectively to the
weakly increasing list of child weights of a root of weight $n$, namely
$k_2$ copies of $2$, then $k_3$ copies of $3$, and so on; the product
$C_2^{k_2}\cdots C_{n-1}^{k_{n-1}}$ is the ordered product of $C_{j_i}$
over this list, and the coefficient $-1/\prod_jk_j!$ is exactly the root's
contribution $(-1)\prod_j1/m_{\mathrm{root},j}!$ to $\omega$. Substituting
the induction hypothesis $C_{j_i}=\sum_{T_i\in\cT_{j_i}}\omega(T_i)A(T_i)$
for each factor and expanding the product gives a sum over all tuples
$(T_1,\ldots,T_r)$ of subtrees, one in each child slot, of
$(-1)\prod_j\frac1{m_{\mathrm{root},j}!}\prod_i\omega(T_i)\prod_iA(T_i)$.
Attaching $T_1,\ldots,T_r$ in order to the root produces exactly one tree
$T\in\cT_n$ with at least one internal vertex; conversely every such tree
arises exactly once (remove the root). For that tree,
$\omega(T)=(-1)\prod_j\frac1{m_{\mathrm{root},j}!}\prod_i\omega(T_i)$
(the internal vertices of $T$ are the root and those of the $T_i$), and
$A(T)=\prod_iA(T_i)$ because the leaves of $T$ in planar order are the
leaves of $T_1$, then of $T_2$, and so on. The remaining term $A_n$ is the
contribution of the single-vertex tree. This proves \eqref{eq:treeformula}.
Finally, $C_n$ is a homogeneous Lie polynomial of degree $n$ by
\cref{thm:zassenhaus}, hence primitive, so \cref{lem:dsw} gives
$C_n=\frac1nD(C_n)=\frac1n\sum_w\coeff{w}C_n\,R(w)$; substituting the word
expansion of \eqref{eq:treeformula} gives \eqref{eq:treelie}.
\end{proof}
For orientation, the first instances are
\begin{equation}\label{eq:CinA}
\begin{aligned}
 C_2&=A_2,\qquad C_3=A_3,\qquad C_4=A_4-\tfrac12A_2^2,\qquad
 C_5=A_5-A_2A_3,\\
 C_6&=A_6-A_2A_4-\tfrac12A_3^2+\tfrac13A_2^3,\qquad
 C_7=A_7-A_2A_5-A_3A_4+\tfrac12A_2^2A_3+\tfrac12A_3A_2^2.
\end{aligned}
\end{equation}
The coefficient $\tfrac13$ of $A_2^3$ in $C_6$ is
$\tfrac12-\tfrac16$: the tree with root children $(2,4)$ and the child of
weight $4$ split into $(2,2)$ contributes $(-1)^2\cdot\frac1{2!}=\frac12$,
and the tree with root children $(2,2,2)$ contributes $-\frac1{3!}$. In
$C_7$ the two terms $A_2^2A_3$ and $A_3A_2^2$ are different polynomials:
the first collects the tree with root list $(2,5)$ and $5\to(2,3)$
(weight $+1$) and the tree with root list $(2,2,3)$ (weight $-\tfrac1{2!}$),
while the second comes from the root list $(3,4)$ with $4\to(2,2)$
(weight $+\tfrac1{2!}$). This illustrates why an unordered product would
lose information. The data file \texttt{zassenhaus\_tree\_polynomials.json}
lists these polynomials in the $A_j$ through degree ten, and the
verification program of \cref{sec:computation} confirms that
\eqref{eq:treeformula} agrees with the two recursive constructions through
that degree.

\subsection{A globally convergent ordered-simplex expansion}
There is another useful all-terms expansion, a sum of ordered products
rather than a product of exponentials; its convergence is therefore a
separate and easier question than Zassenhaus convergence.
\begin{theorem}[Ordered-simplex expansion]\label{thm:simplex}
For bounded $X,Y$ in a unital Banach algebra and every scalar $t$,
\begin{equation}\label{eq:simplex}
\begin{split}
 e^{t(X+Y)}=e^{tX}\Biggl[I+\sum_{k\geq1}\ \sum_{n_1,\ldots,n_k\geq0}
 &\frac{(-1)^{n_1+\cdots+n_k}\,t^{k+n_1+\cdots+n_k}}
 {\displaystyle\prod_{j=1}^{k}n_j!\ \prod_{j=1}^{k}\bigl(n_j+n_{j+1}+\cdots+n_k+k-j+1\bigr)}\\[-4pt]
 &\qquad\cdot\bigl(\ad_X^{n_1}Y\bigr)\bigl(\ad_X^{n_2}Y\bigr)\cdots\bigl(\ad_X^{n_k}Y\bigr)\Biggr],
\end{split}
\end{equation}
and the multiple series is absolutely convergent, with the sum of the norms
of its terms bounded by $\exp\bigl(|t|\norm Ye^{2|t|\norm X}\bigr)$.
\end{theorem}
\begin{proof}
Put $H(t)=e^{-tX}e^{t(X+Y)}$. Differentiating,
$H'(t)=-Xe^{-tX}e^{t(X+Y)}+e^{-tX}(X+Y)e^{t(X+Y)}=e^{-tX}Ye^{t(X+Y)}
=\bigl(e^{-tX}Ye^{tX}\bigr)H(t)=B(t)H(t)$ with $B(t)=e^{-t\ad_X}Y$ by
\cref{thm:campbell}, and $H(0)=I$. First take real $t\geq0$. Integrating,
$H(t)=I+\int_0^tB(s)H(s)\dd s$, and substituting this expression for $H(s)$
repeatedly $N$ times gives
\begin{align*}
 H(t)=I&+\sum_{k=1}^{N}\int_{0<t_k<\cdots<t_1<t}B(t_1)\cdots B(t_k)\,\dd t_k\cdots\dd t_1\\
 &+\int_{0<t_{N+1}<\cdots<t_1<t}B(t_1)\cdots B(t_{N+1})H(t_{N+1})\,\dd t_{N+1}\cdots\dd t_1.
\end{align*}
With $b=\sup_{0\leq s\leq t}\norm{B(s)}\leq\norm Ye^{2t\norm X}$ and
$\norm{H(s)}\leq\norm I^2e^{s\norm X+s\norm{X+Y}}$ bounded on $[0,t]$, the
remainder has norm at most $b^{N+1}t^{N+1}/(N+1)!$ times that bound, since
the ordered simplex has volume $t^{N+1}/(N+1)!$. Hence the remainder tends
to zero and $H(t)=I+\sum_{k\geq1}\int_{\text{simplex}}B(t_1)\cdots B(t_k)$.
Now expand $B(s)=\sum_{n\geq0}(-s)^n\ad_X^nY/n!$ in each factor. The
resulting multiple series of monomials is absolutely convergent: the sum of
the norms of the terms of $B(s)$ is at most
$\sum_n s^n(2\norm X)^n\norm Y/n!=\norm Ye^{2s\norm X}\leq b$, so the sum
of the norms of all monomials in the $k$th ordered integral is at most
$b^kt^k/k!$, and summing over $k$ gives the bound $e^{bt}$. Absolute
convergence justifies integrating term by term. The iterated monomial
integral is evaluated from the inside out: integrating $t_k^{n_k}$ from $0$
to $t_{k-1}$ gives $t_{k-1}^{n_k+1}/(n_k+1)$; then
$t_{k-1}^{n_{k-1}+n_k+1}$ integrates to $t_{k-2}^{n_{k-1}+n_k+2}/(n_{k-1}+n_k+2)$;
continuing,
\[
 \int_{0<t_k<\cdots<t_1<t}\prod_{j=1}^kt_j^{n_j}\,\dd t_k\cdots\dd t_1
 =\frac{t^{k+n_1+\cdots+n_k}}{\prod_{j=1}^{k}(n_j+n_{j+1}+\cdots+n_k+k-j+1)}.
\]
This proves the coefficients in \eqref{eq:simplex} for real $t\geq0$.
For complex $t$, parametrize the segment from $0$ to $t$ by $s=tu$,
$0\leq u\leq1$; the same Picard iteration and the same integrals, now with
the factor $t^{k+\sum n_j}$ pulled out, give the identical formula and the
bound with $|t|$. Multiplying by $e^{tX}$ gives \eqref{eq:simplex}.
\end{proof}

\subsection{An elementary sufficient convergence condition for the product}
Formal stabilization of a product does not automatically imply analytic
convergence. The following bound supplies a concrete domain; it is
deliberately elementary and is not asserted to be optimal.
\begin{theorem}[Convergence of the ordered Zassenhaus product]\label{thm:zconv}
Let $X,Y$ be elements of a unital Banach algebra with $\norm I=1$ and put
$a=\norm X+\norm Y$. Let $r>0$ satisfy
\begin{equation}\label{eq:zconv}
 h(r):=e^{2ar}-1-2ar<2\log2-1.
\end{equation}
Then $\sum_{n\geq2}r^n\norm{C_n}<\infty$, the ordered product in
\eqref{eq:zassenhaus} converges in norm uniformly for $|t|\leq r$, and the
identity \eqref{eq:zassenhaus} holds there. Condition \eqref{eq:zconv} is
satisfied whenever $ar<u_*/2=0.38349\ldots$, where $u_*=0.76699\ldots$ is
the positive solution of $e^{u}-1-u=2\log2-1$; a particularly simple
sufficient condition is $ar\leq\tfrac14$.
\end{theorem}
\begin{proof}
\emph{Majorant.} Expanding the three exponentials in \eqref{eq:R1}, the
sum of the norms of all terms of degree $n$ in $t$ is at most the
coefficient of $t^n$ in $e^{t\norm Y}e^{t\norm X}e^{t(\norm X+\norm Y)}
=e^{2at}$, so $\norm{A_n}\leq(2a)^n/n!$ for $n\geq2$, and
$h(z)=\sum_{n\geq2}(2a)^nz^n/n!$ is a nonnegative-coefficient majorant of
$R_1-I$. Define a scalar formal series $c(z)=\sum_{n\geq2}c_nz^n$ by
\begin{equation}\label{eq:zmajorant}
 c=h+\bigl(e^{c}-1-c\bigr).
\end{equation}
Since $c$ has order $2$, the expression $e^c-1-c=\sum_{m\geq2}c^m/m!$ has
its degree-$n$ coefficient determined by $c_2,\ldots,c_{n-2}$; so
\eqref{eq:zmajorant} determines the $c_n$ recursively, and they are
nonnegative because all ingredients are. In the commutative scalar algebra,
$\prod_{j\geq2}e^{c_jz^j}=e^{c(z)}$, whose coefficient of $z^n$ is
$\sum_{\sum jk_j=n}\prod_jc_j^{k_j}/k_j!$; the choice $k_n=1$ contributes
$c_n$, and the remaining choices, all with at least two factors, contribute
exactly $\coeff{z^n}(e^c-1-c)$. Comparing with \eqref{eq:Cpartitions} and
using the triangle inequality and submultiplicativity,
\[
 \norm{C_n}\leq\norm{A_n}+\sum_{\substack{\sum jk_j=n\\k_n=0}}
 \prod_j\frac{\norm{C_j}^{k_j}}{k_j!}
 \leq h_n+\sum_{\substack{\sum jk_j=n\\k_n=0}}\prod_j\frac{c_j^{k_j}}{k_j!}=c_n,
\]
where the middle inequality uses the induction hypothesis
$\norm{C_j}\leq c_j$ for $j<n$ and monotonicity in the $c_j\geq0$.

\emph{Convergence of the majorant.} The function $q(d)=2d-e^d+1$ is
strictly increasing on $[0,\log2)$, since $q'(d)=2-e^d>0$ there, from
$q(0)=0$ to $q(\log2)=2\log2-1$. Since $h$ is increasing and continuous and
\eqref{eq:zconv} is strict, choose $R>r$ with $h(R)<2\log2-1$, and then
$d\in[0,\log2)$ with $q(d)=h(R)$. Let $\mathcal B$ be the Banach space of
functions holomorphic on $|z|<R$ and continuous on $|z|\leq R$ with the
supremum norm, and let $\mathcal B_d$ be its closed ball of radius $d$.
For $f\in\mathcal B_d$ define $\mathcal F(f)(z)=h(z)+e^{f(z)}-1-f(z)$.
Since $h$ has nonnegative coefficients, $|h(z)|\leq h(R)$ on the closed
disk, and $|e^{f}-1-f|\leq\sum_{m\geq2}|f|^m/m!\leq e^d-1-d$; hence
$\norm{\mathcal F(f)}\leq h(R)+e^d-1-d=q(d)+e^d-1-d=d$, so $\mathcal F$
maps $\mathcal B_d$ into itself. For $f,g\in\mathcal B_d$,
$e^{f}-f-(e^{g}-g)=\int_0^1(e^{g+\tau(f-g)}-1)\,\dd\tau\,(f-g)$, so
$\norm{\mathcal F(f)-\mathcal F(g)}\leq(e^d-1)\norm{f-g}$ with
$e^d-1<1$. By the contraction mapping theorem $\mathcal F$ has a unique
fixed point $f_*\in\mathcal B_d$. Its Taylor coefficients satisfy
\eqref{eq:zmajorant} (coefficient extraction is continuous and
$e^{f_*}-1-f_*$ has the same coefficient recursion), $f_*$ has order at
least $2$ because $h$ does and $e^{f}-1-f$ has order at least twice that of
$f$, and the formal solution of \eqref{eq:zmajorant} is unique; hence
$f_*(z)=\sum_nc_nz^n$ on $|z|<R$. As $c_n\geq0$ and $r<R$,
\[
 \sum_{n\geq2}r^n\norm{C_n}\leq\sum_{n\geq2}c_nr^n=f_*(r)\leq d<\log2.
\]

\emph{Convergence of the product.} Let $P_N(t)=\prod_{n=2}^Ne^{t^nC_n}$.
For $|t|\leq r$ and $M>N$, $P_M-P_N=P_N\bigl(\prod_{n=N+1}^Me^{t^nC_n}-I\bigr)$.
Expanding each factor as $I+(e^{t^nC_n}-I)$ and multiplying out, every term
other than $I$ contains at least one factor $e^{t^nC_n}-I$ of norm at most
$e^{r^n\norm{C_n}}-1$, so
$\norm{\prod_{n=N+1}^Me^{t^nC_n}-I}\leq\exp\bigl(\sum_{n>N}r^n\norm{C_n}\bigr)-1$,
which tends to zero as $N\to\infty$; and $\norm{P_N}\leq\exp(\sum_nr^n\norm{C_n})\leq e^d$.
Hence $(P_N)$ is uniformly Cauchy on $|t|\leq r$ and converges uniformly to
a limit $P(t)$, continuous on the closed disk and analytic on $|t|<r$ as a
uniform limit of entire functions. For $n\leq N$ the coefficient of $t^n$ in
$P_N$ does not depend on $N$ (the omitted factors contribute only degrees
above $N$), and by the Cauchy coefficient formula uniform convergence gives
$\coeff{t^n}P=\coeff{t^n}P_N$. The formal identity of \cref{thm:zassenhaus}
says that $\coeff{t^n}\bigl(e^{tX}e^{tY}P_N(t)\bigr)=\coeff{t^n}e^{t(X+Y)}$
for $n\leq N$. Thus the analytic functions $e^{tX}e^{tY}P(t)$ and
$e^{t(X+Y)}$ have the same Taylor coefficients at $0$, so they agree on
$|t|<r$, and by continuity on $|t|\leq r$.

\emph{The numerical conditions.} The function $u\mapsto e^u-1-u$ is
strictly increasing on $[0,\infty)$ from $0$, so $e^{u}-1-u=2\log2-1
=0.386294\ldots$ has a unique positive solution $u_*$; evaluating,
$e^{0.767}-1-0.767=0.38629\ldots$ and $e^{0.766}-1-0.766=0.38514\ldots$, so
$u_*=0.7669980\ldots$ and \eqref{eq:zconv} holds exactly when
$2ar<u_*$. In particular $ar\leq\tfrac14$ gives
$h(r)\leq e^{1/2}-\tfrac32<1.6488-1.5=0.1488<0.386$.
\end{proof}
The assumption $\norm I=1$ was used only in the bound $\norm{P_N}\leq e^d$;
for a general submultiplicative norm a fixed factor $\norm I$ appears and
nothing else changes.

\subsection{Entire resummation in the eigenvector case}
When $[X,Y]=sY$, the factorization \eqref{eq:eigenfactor} with $c=1$ gives
$e^{t(X+Y)}=e^{tX}e^{q_s(t)Y}$, hence
\[
 R_1(t)=e^{-tY}e^{-tX}e^{tX}e^{q_s(t)Y}=e^{(q_s(t)-t)Y},\qquad
 q_s(t)-t=\sum_{n\geq2}\frac{(-1)^{n-1}s^{n-1}}{n!}t^n.
\]
All exponents are multiples of the same $Y$ and commute, so by
\cref{lem:explog}
\[
 e^{(q_s(t)-t)Y}=\prod_{n\geq2}\exp\Bigl(\frac{(-1)^{n-1}s^{n-1}}{n!}\,t^nY\Bigr),
\]
and uniqueness in \cref{thm:zassenhaus} yields the all-orders closed form
\begin{equation}\label{eq:eigenzass}
 \boxed{\displaystyle C_n=\frac{(-1)^{n-1}s^{n-1}}{n!}\,Y\qquad(n\geq2).}
\end{equation}
This product converges for every finite $s$ and $t$, a useful contrast with
the Bernoulli BCH series \eqref{eq:eigenseries}, whose radius in $s$ is only
$2\pi$. If instead the commutator $C=[X,Y]$ is central, \eqref{eq:centralt}
gives $e^{t(X+Y)}=e^{tX}e^{tY}e^{-t^2C/2}$, so $C_2=-\tfrac12[X,Y]$ and
$C_n=0$ for every $n\geq3$.
